\chapter{Einleitung}\label{ch:einleitung}


\section{Problemstelung und Motivation}

Heutzutage finden \ac{llm}s immer mehr Anwendung in der Praxis und in Unternehmen. Diese Entwicklung wurde auch in vielen Statistiken 
über die Jahre dokumentiert. Dabei stellt das Statistische Bundesamt fest, das in den Jahren 2021 bis 2025 die die Prozentanteile der
Unternehmen welche Künstliche Inteeligenz benutzen von 11 \% auf 26 \% gestiegen. Dabei ist zu beachten das verschiedene Größen von 
Unternehmen miteinbezogen werden, auch welche mit 10-49 beschäftigte. Bei Unternehmen mit einer Beschäftigen Anzahl von 250 und mehr 
stieg die Anzahl von 31 \% auf 57 \%~\cite{e01}.\\
Das IFO Institut verzeichnete, dass im Mai 2026 54,5 \% der Unternehmen in Deutschland 
KI verwenden und 16 \% dies planen. 21,6 \% haben bereits darüber diskutiert und nur 7,9 \% geben an, dass der Einsatz von 
Künstlicher Intelligenz für sie keine rolle spielt~\cite{e02}. 
Der deutlich höhere Anteil gegenüber den Zahlen des Statistischen Bundesamtes ist auf die unterschiedliche Erhebungsmethodik zurückzuführen, 
insbesondere auf eine abweichende Zusammensetzung der befragten Unternehmen und eine weiter gefasste Definition der KI-Nutzung. Beide Erhebungen 
zeigen jedoch übereinstimmend einen deutlichen Anstieg. \\
Dadurch das \acp{llm}s immer häufiger in Unternhemn eingestezt werden greifen sie auch immer mehr auf Dokumete zu oder werden für
weitere Aktionen wie Recherche im Internet eingesetzt. Das wird auch deutlich in dem Bericht der Organisation OWASP von 2026, in dem
die Top 10 Risikien für \ac{llm}s definiert wurden. Dort wurde Excessive Agency auf Platz 3 eingeordnet, wohin gegen 2025 dieses Risiko noch
auf Platz 6 lag. Dieses Risiko umfasst die Gefahr die entsteht, wenn \ac{llm}s auf zu viele Funktionalitäten zugreifen kann und zu viele Rechte zugeteilt bekommt~\cite{pi24}. \\
Dadurch wird deutlich, dass \ac{llm}s immer mehr Zugriff auf Datenmengen gegbene werden, auf welche sie auch 
Aktionen ausführen können. Diese Aktionen werden auch Tool Calls genannt. \\
Mit einem Tool Call werden einem \acp{llm} ermöglicht, dass mehr Aktionen ausgeführt werden können, beispielsweise den Zugriff auf externe Daten
und Webseitem das ausfüllen von Formularen usw.. Es können beliebige Funktionen hinzugefügt werden. Durch diese Funktion entstehen Gefahren, 
da durch den Aufruf externer Quellen und verwendung großer Datenmengen auch Gefahren entstehen 
können. Das geschieht beispielsweise dadurch, das \acp{llm}s nicht zuverlässig zwischen Anweisungen und Daten unterschieden könnn. Somit
können sich in den Daten ebenfalls Anweisungen verstecken, welche dann von dem \acp{llm} ausgeführt werden. Diese Schwachstelle wird von 
Angreifern für sogenannte Prompt Injection angriffe ausgenutzt. Die Angreidfer verstecken Anweisungen in den externn Daten auf die ein
\acp{llm} zugreift. So ist es möglich eschädliche Aktionen auszulösen, ohne direkt Zugriff auf ein \acp{llm} zu haben~\cite{pi15}. \\
Durch die Verwendung von Tool Calls wird dieses Risiko von Prompt Injeciton Angriffen verstärkt, da normalerweise nur eine flasche 
Antwort ausgelöst werden kann, aber dadurch das mit einem Tool Call richtige Aktionen ausgeführt werden können werden 
diese Art von Angriffen immer gefährlicher. Dies wird besonders gefährlich bei Chatbots, welche auf Daten einer Firma zugreifen können, da diese höufig auf viele 
sensible intere Daten zugreifen können und beispielsweise firmengeheimnisse preisgeben können. \\
Dieses Problem ist schon seit mehreren Jahren bekannt, dadurch gibt es auch bereits viele Versuche diese Angriffsfläche einzudämmen.
Es werden teilweise zusätzliche Klassifikatoren, Filter oder Vorverarbeitungen eingesetzt, um zu versuchn die Angriffe bevor sie zu dem
\acp{llm} kommen, welches zugriff auf die Tools hat als ANgriff zu klassifizieren~\cite[4]{pi9}~\cite[9]{sm10}. Jedoch gibt es viele Beweise, dass die Maßnahmen
nicht wirksam genug sind~\cite{sm22}. Die Angreifer verändern ANgriffe immer weiter, damit genau diese umgangen werden.  
Ein bekanntes Beispiel aus dem Jahr 2025 ist der Echoleak von microsoft 365 Copilot, bei dem die Angereifer mithilfe einer vsteckten Anweisung 
in einer E-Mail sensible Unternehmens daten exfiltrieren konnten und sie an einen Server der ANgreifer schickten~\cite{e07}. Für dieses Angriff wurde keine Nutzer
Interaktion notwendig, sondern ausschließlich durch die E-Mail wurde ein Angriff ausgelöst. Dies stellt ein praxisnahes Beispiel für indirekte
Prompt Injeciton dar~\cite{e06}. \\
Manche der Arbeiten, welche in dieser Arbeit berücksichtigt werden, behandeln auch nicht direkt indirekte Prompt Injection Angriffe sondern häufig auch Jailbreaks. 
Der Unterschied ist, dass
Prompt Injection darauf abzielt eine ANwednung welche auf einem \ac{llm} aufbauen, beid dem nicht vertrauenswürdige Eingaben mit einem vertrauenswürdigen Prompt
verkettet werden.
Bei einem Jailbreak hineggebn wird versucht ie Sicherheitsfilter zu umgehen. Die Charakteristik, die Prompt Injection ausmacht ist somit, ob es eine Verkettung von 
einem vertrauenwürdigen und nicht vertrauenswürdigen Prompt gibt~\cite{e05}. \\
Anhand der Literatur wird bewusst das es noch sehr viel Verbesserungspotenzial bei den Schutzmaßnahmen gibt. Genau darum dreht sich diese 
Abreit. Es wird sich auf indirekte Prompt Injectoin Angriffe fokussiert. Bei diesn Angriffen wird über externe Quellen ein Angriff ausgeführt, indem
das Modell auf diese Daten zugreifen soll, zum Beispiel eine Websiten~\cite[2\psq]{pi6}. Es wird versucht eine Schutzmaßnahme zu entwickeln, die 
genauf diese Angriffe abwehrt, welche auch generalisiert angwendet werden kann. Die Schutzmaßnahme wird zunächst mit einem selbst erstellten Datensatz entwickelt, 
dessen Angriffe aus indirekten Angriffsszenarien abgeleitet sind, zur Vereinfachung jedoch direkt im Nutzerprompt übergeben werden. Anschließend wird sie mit dem 
öffentlichen Datensatz BIPIA überprüft, in dem die Angriffe in externe Inhalte wie E-Mails, Tabellen oder Code eingebettet sind.
Als Anwendungsfall wird der Chatvbot ChatBot4U der Firma Energy4U GmbH verwendet, welcher 
auch dazu dienn soll als Unternhemenslösung für die Kommunkation mit Kunden zu dienen, somit ein praxisnaher anwendungsfallbeispiel ist.
Im weiteren wird die genaue Zielsetzung, die Forschungsfrage und der Ablauf geklärt.

\section{Zielsetzung}

Das Ziel dieser Arbeit ist es mit dieser Arbeit herauszufinden, wie stark ein Chatbot, genauer der unternehmensinterne Chatbot von Energy4u GmbH durch eine selbst
entwickelte SChutzmaßnahme vor indirekten Prompt Injection Angriffen abgesichert werden kann. Dafür wird als Vorstufe eine Schwachstellenanalyse durchgeführt,
um die Angriffsfläche zu identfizieren. Für diese Vorarbeit wird eine lokale Testumgebung aufgebaut. \\
Angepasst an diese Schwachstellen wird eine geeignete Schutzmaßnahme entwicklet und implementiert. Es wird anhand mehrerer Datensätze getstet und mithilfe von Metriken
die genaue Leistungsfähigkeit festgehalten. \\
Mit dem Ergebnis der Arbeit soll es möglich sein einen konkreten Vergleich zwischen einer normal implemntierten Schutzmaßnahme und mit einer weiterentwickelten Schutzmaßnahme zu bekommen. Diese
soll sowohl für den unternehmensinternen Chatbot der Energy4U GmbH dienen, sowie übertragbar sein auf andere praxisrelevante Chatbots.

\section{Forschungsfrage}

\begin{enumerate}
    \item HF: Inwieweit kann ein unternehmensinterner Chatvot mit implementierten Tool Calling gegen indirekte Prompt Injections angriffe abgesichert werden?
    \item TF1: Welche Schwachstellen weist der Chatbot überhaupt im Bezug auf indirekte Prompt injection angriffe auf?
    \item TF2: Wie zuverlässig ist ein vorgeschaltetes \ac{llm} in Bezug auf die Abwehr von Prompt injection angriffen?
    \item TF3: Verbessert sich die Erkennung von Prompt Injection Angriffe, wenn die Schutzmaßnahme mithilfe von Segemntierung und Rekonstruktion erweitert wird?
    \item TF4: Inwieweit lassen sich die Ergebnisse auf öffentliche Datensätze übertragen?
\end{enumerate}

Für diese Arbeit wurde ein Forschungsfrage definiert, die wiederrum in mehrere Teilfragen geteilt wurde. Die hauptfrage, welch in dieser Arbeit geklärt werden soll 
ist, inwieweit ein unternehmensinterner Chatvot mit implementierten Tool Calling gegen indirekte Prompt Injections angriffe abgesichert werden kann. Dabei ist der 
Unternhemensinterne Chatvbot nur das Testsystem, die beantwortung der Frage soll aber allgemein gemacht werden können. Zu beantwortung dieser Frage wurden weitere
Fragen definiert. Zum einen soll geklärt werden, welch Schwachstellen der Chatbot überhaupt aufweist im bezug auf indirekte Prompt injection angriffe. Zusätzlich 
sollen bei der Impleneiterung der Schutzmaßnahme die Frage gklärt werden, wie zuverlässig ein vorgeschaltetes \ac{llm} in bzug auf die Abwehr von Prompt injection
angriffe ist. In Bezug auf diese Frage soll zusätzlich geklärt werden, ob sich die Erkennung von Prompt Injection Angriffe verbessert, wenn neben einem vorgeschalteten
Modells, die Schutzmaßnahme mit Rekonstruktionsmaßnahmen und Segmentierungen der Daten erweitert wird. Als letzte soll die Frage beantwortet wredne, inwieweit sich 
die Ergebnisse und der Ansatz auf öffentliche Datensätze übertragen lassen und damit auch wie praxisnah er ist. 

\section{Aufbau der Arbeit}

Diese Arbeit ist in sieben Abschnitte eingteilt. Nach der Einleitung kommt das Grundlagen Kapitel. Dort sollen alle Fachlichen Themen aufarbeitet werden, damit ein
Leser eine erste Übersicht bekommt, worum sich diese Arbeit drehen wird. Dafür werden \acp{llm} behandelt und ihre Funktionsweise. Dazu zählen die Geschihte und 
Technische Bestandteile und ihr Aufbau. Darauf aufbauend wird sich diese Arbeit mit dem Begriff Tool Calling bei \ac{llm}s beschäftigen, also der verwendung von 
externn Ressourcen und das Aufrufen dieser. Das Verständinis über die Funktion ist wichtiger Grundsatz für diese Arbeit. Es werden zudem auch Sicherheitsbedrohnugen,
genannt, die entstehen, wenn man dieses Tool Calling verwendet. Danach werden Chatbots definiert und die ABgrenzung zu \ac{llm}s gemacht. Es wird die Architektur
erläutert und aufgrund diesen technischen Hintergrund die entstehende ANgriffsfläche benannt. Am ende werden noch zusätzlich Unternehmensinterne Chatbots behandelt 
und ihren Einsatzbereich und die Risiken. \\
Nach diesem grundlagen teil, wird das zntrale Thema dieser Arbeit aufgefasst. Es werden Prompt Injeciton Angriffe definiert und eine kategorisierung gemacht in 
indirekte und direkte Angriffe. Zu den beiden Kategorien werden aktuelle Angriffe genannt und erklärt. es werden auch die Bedeutung für Unternehmen definiert und 
ein beispiel genannt, welche als eine überleitung in die praxis dienen soll. \\
Mit diesen Grundlagen können nun der Praxisteil diser Arbeit vorbereitet werden. Dafür wird an dem Unternehmensinternen Chatbot Angriffe getstet.Dfür muss davor 
die Umgeung aufgbaut werden. Es wird zuerst der aufbau des Chatbots erkläuter und wie er sich zusammensetzt. Danach wird die lokale mplemterierung des chatbots, 
sowie alle anderen funktionen und betandteile dieser test umgebung näher erläutert. Anhcsließend wird zu jedem Angriff die durchführung näher gebracht. und die 
erebnisse dokumentiert. Dazu gehören welche Angriffe in welchem Maß erfolgreich waren. Abschließend wurden die dadurch erfassten schwachstellen deieses chatbots
erläutert, die aus den ergebnissen abgeleitet werden könnnen. \\
Da nun die Vorarbeit für den Praxisteil abgeschlossen ist kann nun der hauptaspekt dieser arbeit durchgeführt werden. Es wird eine eigene schutmaßnahme definiert
und implementiert, mit der diese zuvor dokumentierten Angriffe abgewehrt werden sollen. Dafür werden zuerst theoretisch mögliche maßnahme aufbereitet und anahnd 
dueser dokumentierten mit ausgewählten ein konzept entwickelt. Für diese Maßnahme wird ein Datensatz aufgebaut und Metriken ausgewählt. Anschließend wird sie implementiert.
Die gesammelten Ergebnisse werden aufbereitet und für die Auswertung zusammengtragen. \\
Bei der Auswertung wir dzuerst geguckt ob die Fruschungsfrage geklärt werden konnte. Es werden dann die gesammelten Ergebnisse evaluiert und geguckt, ob das
Konzept der Schutzmaßnahme erfolgreich war und warum. Zusätzlich werden die Grenzen definiert und verbesserungsmöglichkeiten. \\
Am ende werden die wichtigstn Ergebnisse der Arbeit gesammelt dokumentiert und ein Ausblick getätigt, wie zukündtige Forschung mögliche Lücken schließen kann.